\documentclass[UTF8,12pt,a4paper,fontset=fandol]{ctexart}

% 支持从项目根目录、深度学习目录或本文件所在目录加载共享样式。
\makeatletter
\def\input@path{{../}{../../}}
\makeatother
\usepackage{researchnotes}

\title{大语言模型中的旋转位置编码}
\author{}
\date{}

\begin{document}
\maketitle
\tableofcontents

\section{研究对象、核心问题与阅读路线}
\label{sec:rope-intro}

\subsection{位置为什么值得单独研究}
考虑“学生帮助老师”和“老师帮助学生”。两句话包含相同的词，关系却因顺序改变而改变。对于逐词读取的循环网络，计算发生的先后顺序已经进入状态更新；对于同时比较整段文本的注意力网络，词元表示之间的相似性本身并不直接标明谁在前、谁在后、相隔多远。\keyterm{位置编码（positional encoding）}就是向模型提供这类结构信息的机制。这里的位置通常指分词器输出序列中的词元索引，而不是汉字编号、单词编号、屏幕上的行号或现实时间。同一句话采用不同分词器时，词元距离可能不同，因此位置机制的语义必须与实际输入表示一起理解。

\keyterm{大语言模型（large language model, LLM）}中的\keyterm{旋转位置编码（rotary position embedding, RoPE）}把每个注意力头内的查询、键向量拆成若干二维向量，再按词元位置旋转这些二维向量。关键并非把位置变成一个额外标签，而是安排一个代数结构：两个向量各自按绝对位置旋转以后，它们的内积可以改写为由相对位移控制的双线性形式。原始方法见 RoFormer 论文\cite{su2021}。本文围绕这一构造独立展开基础推导、算例与实现分析；涉及后续扩展时另行说明来源与适用边界。

用钟表可以获得初步直觉：若两根指针具有相同的角速度，同时前进相同时间，它们之间的角度差不变。RoPE 利用的是这种“共同移动抵消”的性质，但词向量不是只有一根指针。每一对坐标对应一根具有不同角速度的指针，而查询和键原有的方向、长度也参与匹配。因此 RoPE 既不是只看距离的查表函数，也不是把词元内容丢掉以后计算两个位置的相似度。后面的推导将把这三个因素分开：内容决定被旋转的向量，位置决定旋转角度，频率决定位置变化在各坐标平面上的快慢。

\subsection{全文路线与结论层次}
本文先解释自注意力缺少何种位置信息，再补齐平面旋转、正交矩阵与复数的必要知识，推导相对位置恒等式。随后讨论多频率结构、可手算的完整注意力例题、Transformer 中的插入位置、两种坐标排列及参考代码。工程部分集中处理键值缓存、批量填充、滑动窗口、数值精度与梯度；扩展部分讨论训练长度之外的分布变化，以及位置插值、改变基数、分频缩放等方案。最后通过概念辨析、验证协议和带解答的练习，把公式、实现与实验结论连接起来。只需理解基础方法的读者可以先读至第\ref{sec:rope-example}节；需要实际修改模型的读者应继续读缓存与验证部分。

本文严格区分三种陈述。第一种是给定矩阵定义后能够证明的代数性质，例如范数保持、相对位移恒等式；第二种是为了构造可计算模型而作出的设计选择，例如几何频率、只旋转查询和键；第三种是必须由训练与评测支持的经验结果，例如某个缩放方案在特定长度上的任务准确率。代数性质可以帮助理解设计，但不能代替经验验证。特别是，公式对任意位置都能求值，不意味着模型在任意长度都能有效利用信息；注意力分数含有相对位移，也不意味着整个语言模型完全不含绝对位置信息。

\subsection{统一记号与约定}
设序列长度为 $L$，词元位置从 $0$ 开始编号；$m$ 表示查询位置，$n$ 表示键位置，均属于 $\{0,\ldots,L-1\}$。本文固定\keyterm{有向相对位移}为 $\delta=n-m$，因此因果注意力中可见的历史位置满足 $\delta\leq 0$。无向距离写为 $|n-m|$，不得把二者混用。隐藏表示为列向量 $\mathbf h_m\in\mathbb R^D$，$D$ 为模型隐藏维数；单个头的查询、键维数为 $d$，旋转维数为偶数 $d_r=2r\leq d$。先讨论 $d_r=d$，再讨论只旋转部分坐标。$r$ 是二维坐标对的个数，索引 $j$ 从 $0$ 到 $r-1$。

未旋转的查询、键和值分别记为 $\mathbf q_m,\mathbf k_n,\mathbf v_n$，旋转后的查询和键记为 $\widetilde{\mathbf q}_m,\widetilde{\mathbf k}_n$。频率写为 $\omega_j$，位置相位为 $m\omega_j$，频率基数为 $b>1$；这样可以避免把模型参数、旋转角与基数都记成 $\theta$。$\mathbf I$ 表示维数由上下文确定的单位矩阵，$\|\cdot\|_2$ 是欧氏范数，$\odot$ 是逐元素乘法，$\mathrm i$ 是满足 $\mathrm i^2=-1$ 的虚数单位。批量大小另记为 $B$，查询头数与键值头数分别记为 $H_q,H_{kv}$。所有三角函数以弧度为单位，所有对数均为自然对数，除非另有说明。

\section{自注意力到底缺少什么位置信息}
\label{sec:rope-attention}

\subsection{内容寻址与缩放点积}
先固定一个注意力层和一个头，并暂时省去偏置。学习得到的投影矩阵 $\mathbf W_Q,\mathbf W_K\in\mathbb R^{d\times D}$ 与 $\mathbf W_V\in\mathbb R^{d_v\times D}$ 把隐藏表示映射为
\begin{equation}
\mathbf q_m=\mathbf W_Q\mathbf h_m,\qquad
\mathbf k_n=\mathbf W_K\mathbf h_n,\qquad
\mathbf v_n=\mathbf W_V\mathbf h_n.
\label{eq:rope-qkv}
\end{equation}
其中 $d_v$ 是值向量维数，常见设计取 $d_v=d$。\keyterm{缩放点积注意力（scaled dot-product attention）}根据查询与键的内积确定匹配分数。给定位置 $m$ 可以读取的键集合 $\mathcal A_m$，其计算为
\begin{equation}
s_{mn}=\frac{\mathbf q_m^{\mathsf T}\mathbf k_n}{\sqrt d},\qquad
a_{mn}=\frac{\exp(s_{mn})}{\sum_{t\in\mathcal A_m}\exp(s_{mt})},\qquad
\mathbf o_m=\sum_{n\in\mathcal A_m}a_{mn}\mathbf v_n.
\label{eq:rope-attention}
\end{equation}
集合 $\mathcal A_m$ 必须非空。因果注意力通常取 $\mathcal A_m=\{0,\ldots,m\}$；更一般的掩码还可以排除填充或其他文档的词元。本文把禁止读取的位置直接从求和中排除，这与在其分数上加 $-\infty$ 的理想化写法等价。基础架构可参见文献\cite{vaswani2017}。

查询可以理解为“当前位置要寻找什么”，键可以理解为“每个位置能够被什么查询匹配”，值则是“匹配后实际取出的内容”。这是一种便于分析的功能区分，不表示每个向量坐标都有固定的人类语义。分母 $\sqrt d$ 是对内积尺度的调节；在分量近似独立、均值为零且方差受控的启发式条件下，它能够缓解维数增加导致的分数方差增大。训练后的分量并不必然独立，所以这一解释不是对所有网络状态的严格概率结论，也不是 RoPE 采用某组频率的依据。

\subsection{置换等变性与因果掩码的区别}
为说明没有位置机制时的限制，考虑全可见注意力，把行向量排列成 $\mathbf Q,\mathbf K,\mathbf V$。若用同一个置换矩阵 $\mathbf P$ 重排输入位置，则查询、键和值都相应变为 $\mathbf P\mathbf Q,\mathbf P\mathbf K,\mathbf P\mathbf V$。按行计算的 softmax 满足
\[
\operatorname{softmax}(\mathbf P\mathbf S\mathbf P^{\mathsf T})
=\mathbf P\operatorname{softmax}(\mathbf S)\mathbf P^{\mathsf T},
\]
因为左右置换分别只改变行的顺序及每行元素的顺序，指数与归一化操作不依赖原有编号。因此注意力输出变为 $\mathbf P\mathbf O$。这称为\keyterm{置换等变性（permutation equivariance）}：输入重排使输出同样重排，而不是输出保持完全不变。若最后再使用与顺序无关的汇聚，才可能形成置换不变的整体函数。

这个论证假设可见关系也随重排保持一致，不能直接套到固定下三角的因果掩码。因果掩码规定只能读取前缀，已经引入方向与边界，任意打乱词元却保留原掩码会改变计算图。因此“没有位置编码的任何 Transformer 都完全不知道顺序”过于绝对。更准确的说法是：单纯的内容内积没有明确的距离变量，而掩码提供的可见关系与可学习的细粒度距离表示并不是同一机制。RoPE 在这种内容寻址之上，进一步使匹配过程能够依赖相对位移。

\subsection{加法位置编码的交叉项}
一种直接做法是给词元嵌入 $\mathbf e_m$ 加上位置向量 $\mathbf p_m$，令输入为 $\mathbf e_m+\mathbf p_m$。若紧接着应用查询、键投影，记 $\mathbf A=\mathbf W_Q^{\mathsf T}\mathbf W_K$，则点积包含
\begin{equation}
(\mathbf e_m+\mathbf p_m)^{\mathsf T}\mathbf A(\mathbf e_n+\mathbf p_n)
=\mathbf e_m^{\mathsf T}\mathbf A\mathbf e_n
+\mathbf e_m^{\mathsf T}\mathbf A\mathbf p_n
+\mathbf p_m^{\mathsf T}\mathbf A\mathbf e_n
+\mathbf p_m^{\mathsf T}\mathbf A\mathbf p_n.
\label{eq:rope-additive}
\end{equation}
其中既有内容与内容的项，也有内容与绝对位置的交叉项。即使 $\mathbf p_m$ 使用正弦、余弦，任意学习矩阵 $\mathbf A$ 也不保证整个式子只通过 $n-m$ 依赖位置。这不说明加法编码无效，而是说明“位置向量有三角恒等式”与“最终注意力严格具有相对位置形式”之间还隔着投影和交叉项。

RoPE 选择在投影之后构造 $\widetilde{\mathbf q}_m=\mathbf R(m)\mathbf q_m$、$\widetilde{\mathbf k}_n=\mathbf R(n)\mathbf k_n$，希望 $\mathbf R(m)^{\mathsf T}\mathbf R(n)$ 可以化成只由位移决定的矩阵。于是位置结构被直接安排在打分算子中，而不需要训练自行消除式\eqref{eq:rope-additive}的绝对位置交叉项。这里的设计目标是一个精确的代数约束，不是对所有可能位置编码方法的优劣排序。

\section{平面旋转：理解 RoPE 的最小数学准备}
\label{sec:rope-rotation}

\subsection{旋转矩阵及其几何含义}
二维列向量 $\mathbf x=(x_0,x_1)^{\mathsf T}$ 逆时针旋转角度 $\varphi$ 后变为
\begin{equation}
\mathbf R(\varphi)\mathbf x=
\begin{pmatrix}\cos\varphi&-\sin\varphi\\
\sin\varphi&\cos\varphi\end{pmatrix}
\begin{pmatrix}x_0\\x_1\end{pmatrix}
=\begin{pmatrix}x_0\cos\varphi-x_1\sin\varphi\\
x_0\sin\varphi+x_1\cos\varphi\end{pmatrix}.
\label{eq:rope-rotation}
\end{equation}
例如 $(1,0)^{\mathsf T}$ 旋转 $\pi/2$ 得到 $(0,1)^{\mathsf T}$。负号的所在位置不是无关紧要的记法：它决定使用逆时针还是顺时针约定。两种约定都可构造相对位置模型，但推导、频率表、坐标排列和预训练权重必须一致。为了消除歧义，本文始终使用式\eqref{eq:rope-rotation}。

\begin{proposition}[旋转的三个恒等式]
对任意实数 $\alpha,\beta$，有
\begin{equation}
\mathbf R(\alpha)^{\mathsf T}=\mathbf R(-\alpha),\qquad
\mathbf R(\alpha)\mathbf R(\beta)=\mathbf R(\alpha+\beta),\qquad
\mathbf R(\alpha)^{\mathsf T}\mathbf R(\alpha)=\mathbf I_2.
\label{eq:rope-group}
\end{equation}
\end{proposition}
\begin{proof}
第一式由余弦为偶函数、正弦为奇函数直接得到。第二式将两个矩阵相乘，左上角为 $\cos\alpha\cos\beta-\sin\alpha\sin\beta=\cos(\alpha+\beta)$，右上角为 $-\cos\alpha\sin\beta-\sin\alpha\cos\beta=-\sin(\alpha+\beta)$，其余两个元素同理由和角公式得到。把第一、第二式合用，得到 $\mathbf R(-\alpha)\mathbf R(\alpha)=\mathbf R(0)=\mathbf I_2$，即第三式。
\end{proof}

满足 $\mathbf U^{\mathsf T}\mathbf U=\mathbf I$ 的实矩阵称为\keyterm{正交矩阵（orthogonal matrix）}。对旋转矩阵，$\|\mathbf R(\alpha)\mathbf x\|_2^2=\mathbf x^{\mathsf T}\mathbf R(\alpha)^{\mathsf T}\mathbf R(\alpha)\mathbf x=\|\mathbf x\|_2^2$，所以旋转保持长度。若两个向量都旋转相同角度，其内积也保持不变；若旋转角不同，则内积一般改变。RoPE 正是利用后一种情形表达位置关系，同时保留单个向量的长度。把“旋转保内积”不加条件地用在不同位置的向量上，会误得“RoPE 不改变注意力”的结论。

\subsection{从绝对旋转得到相对角度}
设 $\mathbf q,\mathbf k\in\mathbb R^2$，位置 $m,n$ 对应角度 $m\omega,n\omega$。利用式\eqref{eq:rope-group}，有
\begin{equation}
(\mathbf R(m\omega)\mathbf q)^{\mathsf T}
(\mathbf R(n\omega)\mathbf k)
=\mathbf q^{\mathsf T}\mathbf R((n-m)\omega)\mathbf k.
\label{eq:rope-2d-relative}
\end{equation}
左边使用两个绝对位置，右边只显式出现它们的差。这并不是近似抵消，而是在精确算术下成立的恒等式。若把两个位置同时加上 $c$，相对角度不变；若只移动其中一个，角度通常改变。位置间隔在此通过旋转作用于内容向量，因此具有相同位移的不同词元对仍可产生完全不同的分数。

还可以从设计约束反推为什么旋转合理。假定 $\mathbf R(0)=\mathbf I$、每个 $\mathbf R(m)$ 正交，并满足组合规律 $\mathbf R(m+n)=\mathbf R(m)\mathbf R(n)$，那么必有 $\mathbf R(m)^{\mathsf T}\mathbf R(n)=\mathbf R(n-m)$。所以相对位置结构来自“正交性加组合规律”，二维旋转提供了最容易实现的一类构造。本文不声称在不加其他假设时 RoPE 是唯一解；例如所有位置都取单位矩阵也满足相同恒等式，却无法提供有用的位置差异。数学条件限定的是结构，频率与坐标子空间决定结构是否具有足够表达力。

\subsection{方向、面积项与一个符号检查}
把式\eqref{eq:rope-2d-relative}展开，并记 $\delta=n-m$，得到
\begin{equation}
\mathbf q^{\mathsf T}\mathbf R(\delta\omega)\mathbf k
=(q_0k_0+q_1k_1)\cos(\delta\omega)
+(q_1k_0-q_0k_1)\sin(\delta\omega).
\label{eq:rope-sin-cos}
\end{equation}
余弦项对位移反号保持不变，正弦项反号，因此模型一般能够区分前后方向。若这对内容向量恰好平行，正弦项系数为零，则这一坐标平面上的得分对方向不敏感；方向辨别并不是对每个向量对都强制发生。不同平面、不同头以及因果掩码仍可能提供其他方向信息。

检查符号的简便方法是取 $\mathbf q=(1,0)^{\mathsf T}$、$\mathbf k=(0,1)^{\mathsf T}$。当 $\delta\omega=\pi/2$ 时，键旋转为 $(-1,0)^{\mathsf T}$，点积应为 $-1$；式\eqref{eq:rope-sin-cos}也给出 $-\sin(\pi/2)=-1$。若程序或推导得到 $+1$，应先检查位移定义和旋转方向，而不是立即判断两种文献公式冲突。将 $m-n$ 改记为位移时，正弦项的符号必须同时改变。

\section{高维 RoPE 的定义与核心定理}
\label{sec:rope-definition}

\subsection{把一个头拆成多个二维平面}
先设 $d=2r$。把向量按相邻坐标分块为 $\mathbf q_m^{(j)}=(q_{m,2j},q_{m,2j+1})^{\mathsf T}$，键同理。给每个二维平面指定频率 $\omega_j>0$，定义
\begin{equation}
\mathcal R(m)=\operatorname{diag}\bigl(
\mathbf R(m\omega_0),\ldots,\mathbf R(m\omega_{r-1})\bigr),\qquad
\widetilde{\mathbf q}_m=\mathcal R(m)\mathbf q_m,\quad
\widetilde{\mathbf k}_n=\mathcal R(n)\mathbf k_n.
\label{eq:rope-block}
\end{equation}
这里 $\operatorname{diag}$ 表示块对角矩阵：每个二维平面独立旋转，不同平面之间没有由 RoPE 本身引入的混合。注意力投影矩阵则可以在输入特征间学习任意允许的线性组合，所以“旋转时不混合平面”不等于“网络中的信息永远相互隔离”。

\begin{definition}[本文采用的标准频率]
给定基数 $b>1$、旋转维数 $d_r=2r$，定义
\begin{equation}
\omega_j=b^{-2j/d_r},\qquad j=0,1,\ldots,r-1.
\label{eq:rope-frequency}
\end{equation}
对查询与键使用同一组频率和同一种坐标配对。基础 RoPE 中这些频率是预先指定的超参数函数，不是额外的位置嵌入表。
\end{definition}
例如 $b=10000$ 是便于讨论的一种选择，不是任何模型都必须遵守的常数。本文所有涉及该数值的表格均为教学算例，不代表某个未注明版本的实际模型配置。频率也可以按层或按头设计，但每一对参与点积的查询、键必须使用相容的变换，才能保留后面的相对位置证明。

\begin{theorem}[相对位置恒等式与范数保持]
固定频率、坐标配对及内容向量。对任意整数 $m,n,c$，由式\eqref{eq:rope-block}定义的变换满足
\begin{align}
\widetilde{\mathbf q}_m^{\mathsf T}\widetilde{\mathbf k}_n
&=\mathbf q_m^{\mathsf T}\mathcal R(n-m)\mathbf k_n,
\label{eq:rope-main}\\
(\mathcal R(m+c)\mathbf q_m)^{\mathsf T}
(\mathcal R(n+c)\mathbf k_n)
&=\mathbf q_m^{\mathsf T}\mathcal R(n-m)\mathbf k_n,
\label{eq:rope-shift}\\
\|\mathcal R(m)\mathbf q_m\|_2&=\|\mathbf q_m\|_2.
\label{eq:rope-norm}
\end{align}
这些式子同样适用于实数位置，只要变换仍按相位 $m\omega_j$ 定义。
\end{theorem}
\begin{proof}
逐个二维块应用式\eqref{eq:rope-group}，得到 $\mathcal R(m)^{\mathsf T}\mathcal R(n)=\mathcal R(n-m)$。将旋转后的查询、键代入内积即得第一式；将 $m,n$ 同时换成 $m+c,n+c$，两者之差不变，得到第二式。每个二维块都正交，块对角矩阵也正交，从而第三式由二次型计算得到。证明没有使用式\eqref{eq:rope-frequency}的具体形式，因此几何频率影响频谱设计，却不是相对位置恒等式成立的必要条件。
\end{proof}

\subsection{“只依赖相对位置”的准确含义}
式\eqref{eq:rope-main}的显式位置变换只依赖 $n-m$，但 $\mathbf q_m,\mathbf k_n$ 来自隐藏表示，而隐藏表示可能已经汇集了前文、起始符和其他位置信息。因此不能从该式推出 $s_{mn}$ 是一个固定的距离函数 $f(n-m)$。更准确的记法是 $g(\mathbf q_m,\mathbf k_n,n-m)$：内容与位移共同决定得分。同一对词元出现在两个不同上下文中，投影后的内容可能不同，相同相对位移也不保证相同分数。

式\eqref{eq:rope-shift}还具有明确的比较前提：内容向量、可见关系、频率和所有其他运算都保持不变，只改变坐标原点。若在句首插入十个真实词元，前缀集合随之改变，上层隐藏状态也会改变，这不是单纯的坐标平移。若模型有额外绝对位置嵌入、依赖绝对位置的掩码或动态频率切换，同样不能只凭这个局部恒等式断言整体输出保持不变。后面的缓存重编号将专门利用“改变坐标而不改变内容”的情形。

\subsection{部分旋转、奇数头维数与分数尺度}
当 $d_r<d$ 时，分解 $\mathbf q_m=(\mathbf q_{m,R},\mathbf q_{m,U})$、$\mathbf k_n=(\mathbf k_{n,R},\mathbf k_{n,U})$，其中前一部分长度为偶数 $d_r$，后一部分不旋转。于是
\begin{equation}
\widetilde{\mathbf q}_m^{\mathsf T}\widetilde{\mathbf k}_n
=\mathbf q_{m,R}^{\mathsf T}\mathcal R(n-m)\mathbf k_{n,R}
+\mathbf q_{m,U}^{\mathsf T}\mathbf k_{n,U}.
\label{eq:rope-partial}
\end{equation}
未旋转部分提供不显式含位移的匹配通道。偶数约束只针对旋转维数，整个头的维数可以是奇数，只要明确留下哪些坐标不旋转。本文参考实现只旋转前 $d_r$ 个坐标，并按 $d_r$ 构造频率；某个既有检查点若采用其他约定，应以其训练配置为准。

缩放点积的分母仍通常使用实际参与点积的头维数 $\sqrt d$，不能因为旋转了一半坐标就自行改成 $\sqrt{d_r}$。RoPE 保持范数，但改变两个不同位置向量的夹角，因此改变分数并不矛盾。由柯西--施瓦茨不等式还得到 $|\widetilde{\mathbf q}_m^{\mathsf T}\widetilde{\mathbf k}_n|\leq\|\mathbf q_m\|_2\|\mathbf k_n\|_2$。对固定内容向量，这个界不随位置增大而发散；长上下文失效不能简单归因于旋转矩阵把向量长度无限放大。

\section{复数形式、坐标排列与实现等价性}
\label{sec:rope-complex}

\subsection{一个二维实向量就是一个复数}
把 $(x_0,x_1)$ 对应到 $z=x_0+\mathrm i x_1$。欧拉公式给出 $\exp(\mathrm i\varphi)=\cos\varphi+\mathrm i\sin\varphi$，所以
\[
z\exp(\mathrm i\varphi)
=(x_0\cos\varphi-x_1\sin\varphi)
+\mathrm i(x_0\sin\varphi+x_1\cos\varphi).
\]
实部和虚部分别与式\eqref{eq:rope-rotation}的两个坐标相同。复数表达没有引入一种不同的网络，只是把一次二维实旋转写成一次复乘法。实现中采用复数张量时，必须保证每个复数的实部、虚部对应的是训练时约定的那一对实坐标。

设查询和键在第 $j$ 个平面对应 $z_{q,j}=q_{2j}+\mathrm i q_{2j+1}$、$z_{k,j}=k_{2j}+\mathrm i k_{2j+1}$。实内积满足 $\mathbf q^{\mathsf T}\mathbf k=\operatorname{Re}\sum_j\overline{z_{q,j}}z_{k,j}$，其中横线表示复共轭。因此旋转后
\begin{equation}
\widetilde{\mathbf q}_m^{\mathsf T}\widetilde{\mathbf k}_n
=\operatorname{Re}\sum_{j=0}^{r-1}
\overline{z_{q_m,j}}z_{k_n,j}\exp\bigl(\mathrm i(n-m)\omega_j\bigr).
\label{eq:rope-complex}
\end{equation}
若改写为 $z_{q_m,j}\overline{z_{k_n,j}}\exp(\mathrm i(m-n)\omega_j)$，整体是前一个复数的共轭，实部仍相同。这解释了不同材料中位移正负号不同却可能完全等价的原因。真正需要核对的是共轭放在哪一侧、位移采用哪个方向，以及实旋转的正弦符号是否配套。

复数形式还说明 RoPE 是一种内容调制的谐波展开。每个频率前的复系数 $\overline{z_{q_m,j}}z_{k_n,j}$ 由当前查询和键共同决定，既有幅度也有相位；它不是一个全局固定的频率权重。因而不能把 RoPE 简化成“给每个距离乘上一个固定衰减系数”。模型可以通过查询、键投影改变不同频段的使用方式，而频率表本身通常不需要参与梯度更新。

\subsection{相邻配对与前后半段配对}
相邻配对把一个四维向量解释为 $(x_0,x_1)$ 和 $(x_2,x_3)$ 两个平面。另一种常见实现先把坐标排列为 $(u_0,u_1,v_0,v_1)$，再把 $(u_0,v_0)$、$(u_1,v_1)$ 配成两个平面；这称为前后半段配对。若原来的数据是 $(x_0,x_1,x_2,x_3)$，对应的半段排列应为 $(x_0,x_2,x_1,x_3)$，不能在不改变坐标含义的情况下直接换用另一个旋转函数。

用置换矩阵 $\mathbf P$ 表示“相邻排列变为半段排列”，那么两个旋转矩阵满足
\begin{equation}
\mathcal R_{\mathrm{half}}(m)
=\mathbf P\mathcal R_{\mathrm{pair}}(m)\mathbf P^{\mathsf T}.
\label{eq:rope-layout}
\end{equation}
若同时令 $\mathbf q_{\mathrm{half}}=\mathbf P\mathbf q_{\mathrm{pair}}$、$\mathbf k_{\mathrm{half}}=\mathbf P\mathbf k_{\mathrm{pair}}$，两种实现的旋转后内积一致，因为 $\mathbf P^{\mathsf T}\mathbf P=\mathbf I$。这证明它们是同一结构的不同坐标表示，却不意味着可以只替换运行函数而保持投影权重完全不动。转换检查点时，要检查查询、键投影的输出坐标及任何相关的逐坐标参数；只在纸上说明“相差一个置换”还不构成正确转换。

对于相邻配对，定义 $J_{\mathrm{pair}}\mathbf x=(-x_1,x_0,-x_3,x_2,\ldots)$；对于半段配对，若 $\mathbf x=(\mathbf u,\mathbf v)$，定义 $J_{\mathrm{half}}\mathbf x=(-\mathbf v,\mathbf u)$。两者都能写成 $\mathbf x\odot\mathbf c+(J\mathbf x)\odot\mathbf s$，但相邻排列的余弦、正弦表应逐个频率重复两次，半段排列则应把整段频率表复制两次。代码中看见名为 \texttt{rotate\_half} 的函数时，函数名本身不足以判断对应哪一种论文记法。

\subsection{参数与表示之间的边界}
如果从零训练，理论上可以在某个固定置换下同步重参数化查询、键投影，得到等价的函数族。若是使用已经训练的检查点，坐标顺序已经成为参数含义的一部分，随意改变配对会改变模型。这个区别类似于把图像通道从 RGB 改成 BGR：两种表示都合理，输入约定与权重约定却必须一致。判断实现兼容性时，应将坐标排列、频率索引、旋转维数、正弦符号和位置编号一起作为接口契约检查。

\section{多频率结构能表达什么，不能保证什么}
\label{sec:rope-spectrum}

\subsection{频率、波长与相邻位置的变化}
频率 $\omega_j$ 表示位置增加一时角度增加多少弧度，对应的连续位置波长为 $\lambda_j=2\pi/\omega_j$。这里“波长”是完成一整圈需要的连续位置长度；由于词元位置取整数，$\lambda_j$ 不一定是整数步上的精确周期。以 $d_r=8,b=10000$ 为例，四个频率依次为 $1,0.1,0.01,0.001$，波长约为 $6.283,62.832,628.319,6283.185$ 个位置单位。这些量由公式计算得到，反映同一段文本会同时以多个位置尺度进入注意力。

高频平面在相邻位置之间转动较大，更容易产生明显的局部角度差；低频平面转动缓慢，能够在较长范围内保留缓慢变化的相位。这个解释描述的是编码结构，不是对网络内部语义的直接观察。某个频段是否实际承担局部句法或远距离检索，仍取决于训练如何分配查询、键能量。低频坐标若几乎恒为零，即使频率表包含极长波长，模型也未必真正利用了这种尺度。

对固定二维向量 $\mathbf x$，位置相差 $\Delta$ 时，两次旋转结果的距离满足
\begin{equation}
\|\mathbf R((m+\Delta)\omega)\mathbf x-
\mathbf R(m\omega)\mathbf x\|_2
=2\|\mathbf x\|_2\left|\sin\frac{\Delta\omega}{2}\right|.
\label{eq:rope-chord}
\end{equation}
证明只需把平方展开，利用同角旋转保内积，得到 $2\|\mathbf x\|_2^2(1-\cos(\Delta\omega))$，再使用半角公式。小角度时右边约为 $\|\mathbf x\|_2|\Delta\omega|$。因此低频通道中相邻位置的差别较小，但这只是可分辨性的一个局部尺度分析，不能直接推出任务错误率。

\subsection{为什么采用几何频率而不是一个频率}
单一频率容易反复出现近似相位，远处两个位置可能在这个平面上非常接近。多频率把位置映射到多个圆的乘积中，各平面的近似重合通常不会在同一距离同时发生，从而提供更丰富的区分信息。几何序列的相邻频率比固定，即 $\omega_{j+1}/\omega_j=b^{-2/d_r}$，相当于在对数尺度上均匀布置频率。这是一种覆盖多个尺度的简洁设计，并不存在一个从一般语言规律唯一推出 $b$ 的定理。

增加基数 $b$ 时，$\omega_0=1$ 保持不变，其他 $\omega_j$ 随 $j>0$ 变小；所以改变基数不是把所有频率统一除以同一个数。旋转维数固定时，这会拉宽波长范围，同时改变频率之间的间距。增加旋转维数则在同一公式下增加频率采样点，并使最低频率更接近 $1/b$，但它还牵涉头维数、模型容量和计算预算。讨论“更大的基数更好”或“更多频率更好”时，必须同时给出训练长度、数据、优化和评价协议。

将 $b$ 称为“最大上下文长度”也不准确。以式\eqref{eq:rope-frequency}为准，最慢频率是 $b^{-(d_r-2)/d_r}$，对应波长为 $2\pi b^{(d_r-2)/d_r}$，它既不等于 $b$，也不是系统必须停止生成的位置。模型允许的窗口由训练配置、位置策略、掩码、缓存容量和运行时实现共同决定，不能从频率表中的一个数字直接读出。

\subsection{周期性、混叠与近似碰撞}
\keyterm{混叠（aliasing）}在这里指不同位置或位移经周期函数映射后出现相同或接近的表示。若全部频率都是某个基本频率的整数倍，并且存在整数 $T$ 使 $T\omega_j\in2\pi\mathbb Z$，则 $\mathcal R(m+T)=\mathcal R(m)$ 对所有平面成立。几何频率并不一般满足具有小整数 $T$ 的条件，因此不能把单个波长当成完整高维编码的共同周期。

反过来，没有短的精确共同周期也不意味着任意长范围都能稳定区分。有限维表示落在有界空间中，考虑越来越多的位置，总会有某些位置编码相距很近；有限数值精度还会进一步限制区分能力。对于有限训练范围，重要的是任务所需要区分的位移是否具有足够大的有效差异，而不是声称一个有限维周期编码可以无误差地表示无限多种距离。这样的有效差异还受到内容系数的影响，因为模型实际使用的是加权内积，不是直接读取完整相位向量。

\subsection{为什么“远距离衰减”不是逐点定理}
固定一对内容向量，并暂时省去位置下标，RoPE 内积具有
\begin{equation}
F(\delta)=\sum_{j=0}^{r-1}
\left[A_j\cos(\delta\omega_j)+B_j\sin(\delta\omega_j)\right],
\quad
A_j=q_{2j}k_{2j}+q_{2j+1}k_{2j+1},\quad
B_j=q_{2j+1}k_{2j}-q_{2j}k_{2j+1}.
\label{eq:rope-fourier}
\end{equation}
这是有限个振荡项的和。不同频率的项可能在某些距离抵消，也可能在另一些距离重新增强。若只让一个平面非零，并取 $\mathbf q^{(0)}=\mathbf k^{(0)}=(1,0)^{\mathsf T}$，就得到 $F(\delta)=\cos(\delta\omega_0)$；它会反复升降，甚至变号，并不随 $|\delta|$ 单调减小。正交性也不迫使这个分数的绝对值随距离趋于零。

有些分析会观察平均位置相关函数 $K(\delta)=r^{-1}\sum_j\cos(\delta\omega_j)$，或在一定系数规律下估计频率抵消形成的包络。这类对象有助于理解多尺度偏置，但 $K$ 不是任意训练状态下的真实注意力分数，包络也不是每个样本、每个头的单调性保证。若引入随机内容假设，还要说明概率分布；例如相互独立、零均值的查询和键可使内积期望对所有距离都为零，这样的“平均”显然不能证明模型更偏爱近处。任何统计衰减论断都应说明究竟平均了什么。

另外，注意力权重还经过 softmax。即使某个未归一化分数降低，其权重也可能因其他候选分数降低得更多而升高。比较两个不同长度的序列时，分母包含的候选数也不同。因此必须区分旋转内积、缩放后的分数、归一化权重及最终输出四个层次。把位置编码的相关性曲线直接称为“模型的注意力曲线”，会掩盖内容变化与归一化竞争。

\subsection{对位置扰动的一个可核查界}
对固定内容向量，式\eqref{eq:rope-fourier}对实变量 $\delta$ 可微。由于二维旋转的导数为 $\omega_j\mathbf J\mathbf R(\delta\omega_j)$，其中 $\mathbf J=\begin{pmatrix}0&-1\\1&0\end{pmatrix}$ 也是正交矩阵，柯西--施瓦茨不等式给出
\begin{equation}
|F'(\delta)|\leq\sum_j\omega_j
\|\mathbf q^{(j)}\|_2\|\mathbf k^{(j)}\|_2,
\qquad
|F(\delta+\varepsilon)-F(\delta)|
\leq|\varepsilon|\sum_j\omega_j
\|\mathbf q^{(j)}\|_2\|\mathbf k^{(j)}\|_2.
\label{eq:rope-lipschitz}
\end{equation}
第二式由对第一式积分得到。它说明位置误差的影响同时取决于频率和该频段的内容能量；只看基数大小无法完整描述敏感性。这个界可能很松，也不直接约束多层模型的输出误差，但它能解释为什么测试位置编号偏一、精度丢失或频率不一致时，应同时覆盖不同频段和不同内容方向。

\section{一个从旋转到注意力输出的完整算例}
\label{sec:rope-example}

\subsection{给定内容与位置，先算相对内积}
设一个注意力头的维数 $d=4$，采用 $b=10000$，故两个平面的频率是 $1$ 和 $0.01$。当前位置 $m=2$ 的查询为 $\mathbf q_2=(1,0,1,0)^{\mathsf T}$，三个可见位置的键分别为
\[
\mathbf k_0=(1,0,1,0)^{\mathsf T},\qquad
\mathbf k_1=(0,1,1,0)^{\mathsf T},\qquad
\mathbf k_2=(1,0,0,1)^{\mathsf T}.
\]
这些向量是人工指定的教学数据，不来自某个模型的实际中间状态。其目的是把内容与位置的作用分开，并展示每个步骤如何核查。

位置 $n=0$ 对应 $\delta=-2$，两个平面都是相同方向的向量，所以未缩放分数为 $\cos2+\cos0.02$。位置 $n=1$ 对应 $\delta=-1$，第一个平面是查询 $(1,0)$ 与键 $(0,1)$，其贡献为 $-\sin(-1)=\sin1$；第二个平面贡献为 $\cos0.01$。位置 $n=2$ 对应 $\delta=0$，两个平面贡献分别为 $1$ 和 $0$。于是
\begin{equation}
s_{20}=\frac{\cos2+\cos0.02}{2}\approx0.291827,\quad
s_{21}=\frac{\sin1+\cos0.01}{2}\approx0.920710,\quad
s_{22}=\frac12.
\label{eq:rope-example-score}
\end{equation}
分母为 $\sqrt4=2$。尽管位置 $0$ 的键在内容上与查询完全相同，旋转后位置 $1$ 的分数却更高；这展示了位置结构如何改变内容匹配，而不是只在已有注意力权重上附加一个标签。

\subsection{归一化、值向量与结果解释}
令三个值向量为 $\mathbf v_0=(1,0)^{\mathsf T}$、$\mathbf v_1=(0,2)^{\mathsf T}$、$\mathbf v_2=(-1,1)^{\mathsf T}$。把式\eqref{eq:rope-example-score}代入 softmax，得到 $(a_{20},a_{21},a_{22})\approx(0.243490,0.456670,0.299840)$，所以 $\mathbf o_2=(a_{20}-a_{22},\,2a_{21}+a_{22})^{\mathsf T}\approx(-0.056350,1.213180)^{\mathsf T}$。为了数值稳定，实际计算可先从三个分数都减去其中最大值；共同平移不改变归一化概率。注意值向量的维数在这个例子中为二，不妨碍查询和键使用四维旋转，因为值只在权重确定后参与加权求和。

若完全去掉 RoPE，三个内容内积分别为 $2,1,1$，除以二后为 $1,0.5,0.5$，此时第一个位置权重最高。比较两种结果可以看出，RoPE 通过查询与键控制“读哪里”，值向量决定“读出什么”。标准构造中不旋转值并不意味着输出不含位置信息，因为输出使用的权重已经是位置的函数。反之，如果额外旋转值，就会改变加权求和的内容坐标系，那属于需要重新分析的架构修改。

\subsection{同时平移位置与改变真实上下文}
把三个键的位置改记为 $10,11,12$，查询改记为 $12$，同时保持给定查询、键、值向量与可见关系不变，式\eqref{eq:rope-example-score}完全不变。读者可以分别计算绝对旋转，再与相对形式比较，作为实现的最小检查。然而在真实语言模型前面插入十个新词元后，查询、键和值通常不再是这里指定的向量，不能据此预测新输出必然相同。这个算例核验的是旋转层的代数结构，不是完整模型对任意前缀插入的不变性。

同样，把所有位置从零起始改为从一起始，只要对整个序列一致处理，并且没有引入额外绝对位置机制，就属于共同平移。只给查询增加一而不给缓存键增加一，则改变所有查询到历史键的相对位移。这种错误在初始预填充阶段可能不出现，却会在逐词解码时持续影响输出，因此不能仅用一次全序列前向检查缓存实现。

\section{RoPE 在 Transformer 层中的具体位置}
\label{sec:rope-architecture}

\subsection{先投影，再旋转，再计算注意力}
在一种常见的预归一化解码器层中，隐藏状态先经过归一化，再分别投影为查询、键和值；查询和键按位置旋转，随后执行带掩码的缩放点积注意力，拼接各头输出，最后经过输出投影与残差连接。之后的前馈网络不需要为了使用 RoPE 再次施加同一旋转。不同模型的归一化与投影顺序可能有变化，本文只明确 RoPE 的关键接口：它作用于用于打分的查询、键坐标，而不是默认作用于整个隐藏状态。

若把旋转挪到投影之前，原来的 $\mathcal R(m)\mathbf W_Q\mathbf h_m$ 会变成 $\mathbf W_Q\mathcal R_D(m)\mathbf h_m$，其中后一个旋转甚至处在不同维数的空间。这两个式子一般不相等；即便维数相同，任意学习矩阵也不与旋转矩阵交换。只有在专门的交换或重参数化条件下，才可能得到等价实现。因此，优化代码时不能把一个固定位置上的线性变换随意穿过另一个线性变换，也不能为了减少调用次数就把 RoPE 移到词嵌入层。

RoPE 通常在多个注意力层分别作用于该层新产生的查询、键，而不是对上一层已经旋转的查询继续旋转。原因是每层都从更新后的隐藏表示重新构造查询、键；这些向量具有新的内容含义。若程序误把已旋转查询当作未旋转输入，再按同一位置旋转一次，就会得到 $\mathcal R(2m)\mathbf q_m$，相当于改变了相位。每层重复使用同一种位置规则，与同一向量被重复施加位置变换，是两件不同的事。

\subsection{多头与分组查询注意力}
\keyterm{多头注意力（multi-head attention, MHA）}使用不同投影产生多个查询、键和值子空间；即便各头共享相同频率表，不同的内容投影也能形成不同的位置使用模式。频率表相同不表示所有头具有相同的注意力分布。\keyterm{分组查询注意力（grouped-query attention, GQA）}令多个查询头共享一个键值头，\keyterm{多查询注意力（multi-query attention, MQA）}则只保留一个键值头。它们改变的是键值共享关系，而不是式\eqref{eq:rope-main}本身。

设 $H_q$ 是 $H_{kv}$ 的整数倍，每个查询头被映射到一个键值头。与同一共享键进行点积的查询必须采用与这个键相容的旋转规则；若随意给这些查询头设置不同频率，而缓存键只有一种旋转版本，则一般不能同时满足各头的相对位置恒等式。工程上可以先旋转少量键值头，再按分组规则广播键；当复制的各头使用同一频率时，复制与旋转可以交换。广播维度、头映射和旋转维数必须一起核对。

基础实现可以把查询、键张量排列为 $[B,H,L,d]$，也可以排列为 $[B,L,H,d]$。数学公式不偏好某一种布局，但正弦、余弦表的广播维度依赖实际排列。若频率表为 $[B,L,r]$，在第一种布局中通常需要插入长度为一的头轴，得到 $[B,1,L,r]$。许多形状错误并不会立即报错：例如当批量大小与头数碰巧相同时，错误广播仍能运行，所以验证应使用互不相等的维数。

\subsection{为什么通常不旋转值}
RoPE 的核心目标已经通过查询和键的双线性匹配实现。保持 $\mathbf v_n$ 的内容坐标不变，使注意力输出仍是这些值向量在同一坐标系中的加权和。若同时将值变为 $\mathcal R(n)\mathbf v_n$，不同位置的内容会先被置于不同旋转坐标，再直接相加；这可以构成某种其他设计，但需要解释输出投影如何处理它，并重新训练或证明相应变换关系。不能因为查询、键、值都由线性投影产生，就认为三者应完全对称地处理位置。

在某一层不旋转值，也不代表值永远不包含位置线索。值来自隐藏状态，隐藏状态可能已经由前面各层的带位置注意力更新过。这里区分的是“本层是否显式给值施加旋转”与“值是否统计上携带位置信息”。前者是代码可直接检查的结构事实，后者则涉及网络表示，需要更细致的分析或干预实验。

\section{参考实现：让张量操作与公式逐项对应}
\label{sec:rope-code}

\subsection{接口约定与计算过程}
下面给出教学用 PyTorch 实现，输入 \texttt{x} 的形状固定为 $[B,H,L,d]$，\texttt{position\_ids} 为 $[B,L]$，返回形状相同的张量。函数只处理一个查询或键张量，不隐式生成位置，不修改掩码，也不管理缓存。这样，位置编号的责任留在调用者，适合比较预填充与增量解码。代码采用相邻配对，仅旋转前 \texttt{rotary\_dim} 个坐标。为便于查看，基本频率计算与旋转写在同一函数中；实际系统可以缓存固定频率及三角函数表。

位置在转成浮点之前应保存为整数索引；若要实现位置插值，也可以传入明确设计的实数坐标。本函数在普通低精度输入下用 float32 计算相位与旋转，在输入为 float64 时保留 float64，以便建立数值参考。禁止自动混合精度作用于该计算块，是为了避免外层训练上下文改变三角函数前的计算精度。此处的代码强调可读性与数学契约，并未针对特定设备设计融合内核。

\begin{verbatim}
import torch

def rope_pairwise(x, position_ids, base=10000.0,
                  rotary_dim=None):
    if x.ndim != 4 or not x.is_floating_point():
        raise ValueError("x must be floating [B, H, L, D]")
    batch, heads, length, dim = x.shape
    if position_ids.shape != (batch, length):
        raise ValueError("position_ids must be [B, L]")
    dr = dim if rotary_dim is None else rotary_dim
    if not isinstance(dr, int) or dr <= 0 or dr > dim:
        raise ValueError("invalid rotary_dim")
    if dr % 2 != 0 or not base > 1.0:
        raise ValueError("need even rotary_dim and base > 1")
    work = torch.float64 if x.dtype == torch.float64 else torch.float32
    with torch.autocast(device_type=x.device.type, enabled=False):
        ids = position_ids.to(device=x.device, dtype=work)
        j = torch.arange(dr // 2, device=x.device, dtype=work)
        inv_freq = base ** (-2.0 * j / dr)
        angle = ids[..., None] * inv_freq
        c = angle.cos()[:, None, :, :]
        s = angle.sin()[:, None, :, :]
        xr = x[..., :dr].to(work)
        u, v = xr[..., 0::2], xr[..., 1::2]
        y = torch.stack((u * c - v * s,
                         u * s + v * c), dim=-1)
        y = y.flatten(-2).to(x.dtype)
    return torch.cat((y, x[..., dr:]), dim=-1)

# q: [B, Hq, Lq, D], k: [B, Hkv, Lk, D]
# Query and key positions are supplied separately.
# q_rot = rope_pairwise(q, q_position_ids)
# k_rot = rope_pairwise(k, k_position_ids)
\end{verbatim}

代码中的 \texttt{inv\_freq} 实际存储本文的角频率 $\omega_j$。名称来自频率与波长倒数的关系，不表示还需再次求倒数。\texttt{stack} 先恢复每个二维平面的两个坐标，再用 \texttt{flatten} 恢复原维数；如果直接把全部第一坐标与全部第二坐标拼接，就会意外转成半段排列。位置表在头轴上广播，允许同一批次中不同样本具有不同位置编号，也允许查询与键长度不相同。

\subsection{半段实现和复数实现如何交叉核对}
对旋转部分长度 $d_r=2r$ 的半段表示，取前 $r$ 个坐标为 $\mathbf u$、后 $r$ 个坐标为 $\mathbf v$，分别计算 $\mathbf u\odot\cos-\mathbf v\odot\sin$ 和 $\mathbf u\odot\sin+\mathbf v\odot\cos$，再按半段顺序拼接即可。验证两种代码时，应先按式\eqref{eq:rope-layout}置换输入，再把输出逆置换回来比较。直接把相同数组同时送入两种配对函数，并期待输出逐项相同，是一个错误的测试。

复数实现则把最后一维重排为 $[r,2]$，把二元组视为实部与虚部，再乘以 $\exp(\mathrm i m\omega_j)$。Meta 的公开 Llama 参考实现展示了这种计算路径\cite{meta-llama}；本文代码没有逐行复制该实现，而是采用显式实数运算便于检查符号。复数视图常对数据类型与内存步幅有要求，不能把任意非连续张量无条件视为复数。实数版本避免了这类视图约束，但并不能仅凭写法判断哪种版本在某设备上更快。

\subsection{计算量与预计算的收益边界}
直接构造每个位置的 $d_r\times d_r$ 矩阵再进行一般矩阵乘法，会掩盖块对角稀疏结构。逐坐标公式对每个词元只需要与 $d_r$ 成正比的乘加，查询、键的总旋转计算量为 $O(BL(H_q+H_{kv})d_r)$。在稠密自注意力中，成对打分的主要计算量通常随序列长度平方增长，所以 RoPE 不会把一般稠密注意力自动变为线性复杂度。它的额外开销与注意力本身的复杂度应分开计算。

预计算某个长度范围的正弦、余弦表可以减少重复三角函数求值，但需要 $O(Ld_r)$ 的存储，通常可在批次和头之间共享。若不同层使用不同频率、不同请求具有不同缩放配置，则共享范围会缩小；若位置很长而实际只访问少数位置，按需计算也可能更合适。是否缓存完整表、缓存频率或融合旋转应由运行负载决定，不能把某一种代码优化描述为数学定义的一部分。

\section{键值缓存、位置编号与增量解码}
\label{sec:rope-cache}

\subsection{预填充与逐词生成的同一坐标系}
\keyterm{键值缓存（key--value cache, KV cache）}保存历史位置的键和值，使自回归生成无需反复计算相同历史状态。\keyterm{预填充（prefill）}一次处理已有提示词，\keyterm{解码（decoding）}随后逐步处理新增词元。在固定频率的标准实现中，可以直接缓存已经旋转的键 $\widetilde{\mathbf k}_n=\mathcal R(n)\mathbf k_n$ 与未经本层 RoPE 旋转的值 $\mathbf v_n$。新位置 $m$ 只需计算并旋转自己的查询和新键，再与缓存中可见键计算内积。

例如提示长度为 $P$，位置编号是 $0,\ldots,P-1$，下一次处理新词元的位置应为 $P$。如果每次单词元前向都使用本地长度为一的 \texttt{arange}，位置会反复变成零，导致新查询与旧键处于错误的相对坐标关系。位置索引应来自请求的逻辑序列状态，而不能仅从当前张量长度推断。分块预填充同样如此：第二个分块的张量索引可以从零开始，逻辑位置必须接在第一个分块之后。

在固定频率、相同掩码且没有训练随机性干扰的条件下，全序列因果前向与逐步缓存前向应产生一致的数学结果，因为历史位置不能读取未来位置，它们的隐藏状态不会因未来词元补入而改变。实际数值比较需要允许不同矩阵乘法形状造成的舍入差异，并关闭 dropout 等随机操作。比较最终生成文本太粗糙：两种实现可能在若干步恰好选择相同词元，却已经具有明显不同的 logits。

\subsection{填充位置、注意力掩码与打包文档}
批量推理时，各样本长度通常不同，需要\keyterm{填充（padding）}使张量对齐。位置编号与是否可见是两个不同问题：\texttt{position\_ids} 指定真实词元采用哪个逻辑位置，注意力掩码决定哪些键可以被读取。给填充词元分配零号位置并不能阻止真实词元读取它；反过来，掩掉填充也不能自动修正真实词元的位置编号。对于常见的二维有效词元掩码，可以用沿序列轴的累计有效计数减一构造连续位置，再给填充处填入安全占位值，但要确认这与模型预处理约定一致。

左填充和右填充在数组中产生不同索引，理想的纯 RoPE 机制允许真实序列采用一致的逻辑编号。某些实现给整个样本使用共同偏移，理论上也可由平移性质抵消，但额外绝对位置机制、缓存接口和有限精度可能使这种处理不再等价。因此更清楚的做法是明确记录真实位置，再通过测试验证不同填充方式下的真实词元输出。尤其不能把批次内最长样本的历史长度无条件用作所有样本下一词元的位置。

将多篇独立文档打包到一个训练序列时，若希望文档间彼此不可见，应同时设置块状因果掩码，并按所选训练方案处理每篇文档的位置。仅重置位置而不阻断跨文档注意力，会让不同文档中的词元发生非预期交互；仅阻断注意力却使用连续编号，在纯固定频率 RoPE 的理想条件下可与各文档重置原点等价，但应考虑起始符等其他机制。因果性、文档隔离和坐标原点各自承担不同职责。

\subsection{滑动窗口删除内容与重编号}
\keyterm{滑动窗口注意力（sliding-window attention）}只保留最近一段可见历史。删除旧缓存以后，可以继续使用单调增长的全局位置；这样保留键不需要重新旋转。另一种做法是把保留窗口的起点重新记为零，这会改变其坐标原点。若旧位置为 $n$，新位置为 $n-c$，则已旋转键必须变为
\begin{equation}
\widetilde{\mathbf k}'_{n-c}
=\mathcal R(-c)\widetilde{\mathbf k}_n
=\mathcal R(n-c)\mathbf k_n.
\label{eq:rope-rebase}
\end{equation}
查询也必须按同一新原点旋转，才能保持保留位置之间的分数不变。只改变元数据中的位置编号而不调整已旋转缓存键，会使缓存内容与其标注不一致。

这里保持的是固定缓存内容之间的局部打分关系。删除旧词元可能改变模型能够访问的信息，而且现有隐藏状态还可能包含过去读取过的旧内容。将当前缓存重编号，并不等于对截断后窗口从头重新运行整个多层网络。两者的差别来自上下文历史，不是旋转恒等式失效。分析流式模型时，要区分“保留缓存的坐标变换”与“重新计算一个更短输入”的语义。

\subsection{动态频率与缓存一致性}
若生成长度增加后把频率从旧配置改为新配置，旧键中保存的是 $\mathcal R_{\mathrm{old}}(n)\mathbf k_n$，新查询却使用 $\mathcal R_{\mathrm{new}}(m)\mathbf q_m$。两者的点积通常不能化成单一频率系统下的相对位移。这是一个结构不一致问题，即使每个旋转函数单独计算都正确，也可能产生错误结果。它不同于采用固定大基数从头训练或在一个请求开始时固定缩放配置。

在只考虑固定未旋转键的情况下，可以通过
\begin{equation}
\widetilde{\mathbf k}_{n,\mathrm{new}}
=\mathcal R_{\mathrm{new}}(n)
\mathcal R_{\mathrm{old}}(n)^{\mathsf T}
\widetilde{\mathbf k}_{n,\mathrm{old}}
\label{eq:rope-cache-convert}
\end{equation}
把旧键转换到新频率，或缓存未旋转键并重新施加位置变换。不过多层网络还有更深一层限制：历史未旋转键来自在旧频率下形成的隐藏状态。仅转换当前缓存的旋转部分，一般不能恢复“整个前缀始终在新频率下重新前向”的所有隐藏状态和值。前者修复局部表示约定，后者是更强的全模型一致性要求，必要时需要重新计算历史前缀。

实际系统可以在请求开始时固定位置策略，也可以明确实现支持动态切换的缓存协议，但应把协议作为模型语义的一部分验证。缓存的有效性不仅由词元内容决定，还依赖层、坐标排列、旋转维数、位置原点、频率与缩放版本。复用前缀缓存时忽略这些条件，可能得到维数正确却语义错误的注意力结果。

\section{数值精度、反向传播与稳定性边界}
\label{sec:rope-numerics}

\subsection{位置应在什么精度下计算}
数学上的相邻位置 $m$ 和 $m+1$ 不同，但低精度浮点数未必能区分它们。以通常的二进制浮点格式为例，bfloat16 在数值达到 $256$ 附近后便不能表示所有相邻整数，float16 在 $2048$ 附近也出现同类问题；float32 的连续整数精确表示范围大得多，但超过 $2^{24}$ 后同样不能表示每一个整数。因此，若先把长位置索引变为低精度，再计算相位，两个不同位置可能在三角函数之前就已变成同一个数。

先用整数保存位置，再在较高精度下构造相位，能够减少这种风险。需要注意，最后把结果转成 float32 并不能恢复之前已经丢失的整数差异：\texttt{position\_ids} 若先经过 bfloat16 舍入，再升精度，丢掉的信息仍然丢掉。频率表也应从足够精确的原始参数构造，不能反复把已经舍入的低精度表作为高精度基准。为了验证实现，可以用 float64 计算相同样例的参考值，再考察实际设备和目标精度下的误差。

即使位置整数被精确表示，乘积 $m\omega_j$ 仍可能舍入，三角函数也有自身的实现误差。若相位误差为 $\varepsilon$，由式\eqref{eq:rope-chord}可知，同一二维向量旋转输出的误差至多为 $\|\mathbf x\|_2|\varepsilon|$，因为 $2|\sin(\varepsilon/2)|\leq|\varepsilon|$。因此绝对位置很大时，关注的应是最终相位误差，而不仅是输出张量采用了哪种数据类型。对超出通常工作区间的极大位置，需要单独评测相位生成与三角函数精度。

\subsection{共同平移在浮点数中只能近似成立}
在实数精确运算中，位置同时增加一个很大的常数不影响相对内积。然而程序分别计算 $\cos(m\omega_j)$ 和 $\cos(n\omega_j)$ 时，大绝对相位可能具有不同舍入误差，所以共同平移测试应使用合理容差，并覆盖原始训练长度与实际部署长度。一个实现如果在小位置通过测试，在很大的位置偏差明显，应先检查相位计算，再判断是理论设计还是模型能力问题。

把所有位置减去某个共同原点可以改善相位规模，但与缓存同时使用时必须按照式\eqref{eq:rope-rebase}保持一致。若系统需要很长的流式历史，也不能认为不停重置计数器即可无条件解决所有问题：反复重旋转会积累舍入误差，保留内容的语义和窗口策略仍需独立分析。数值策略是实现近似的一部分，应通过目标长度上的误差测量选择，而不是用代数等价来代替数值评估。

\subsection{对查询、键和频率的梯度}
对固定位置与频率，RoPE 是线性变换。若损失为标量 $\mathcal L$，则链式法则给出
\begin{equation}
\frac{\partial\mathcal L}{\partial\mathbf q_m}
=\mathcal R(m)^{\mathsf T}
\frac{\partial\mathcal L}{\partial\widetilde{\mathbf q}_m}.
\label{eq:rope-backward}
\end{equation}
由于矩阵正交，旋转这一局部步骤在精确算术下保持反传梯度的欧氏范数。该结论只针对这个固定线性算子，不能推导出整个 Transformer 永不梯度爆炸或消失；投影、归一化、softmax、残差叠加和训练目标都参与最终梯度行为。

若把频率也设为可学习参数，则第 $j$ 个二维块满足
\[
\frac{\partial}{\partial\omega_j}
\bigl(\mathbf R(m\omega_j)\mathbf q_m^{(j)}\bigr)
=m\mathbf J\mathbf R(m\omega_j)\mathbf q_m^{(j)},
\qquad \mathbf J=\begin{pmatrix}0&-1\\1&0\end{pmatrix}.
\]
频率梯度中出现绝对位置系数 $m$，所以学习频率和学习普通查询权重具有不同的尺度问题。若目标只通过成对相对内积依赖位置，则直接对该内积求频率导数会出现 $n-m$；不同计算图中的中间梯度不应混为一谈。可学习频率属于一种扩展设计，需要额外讨论初始化、约束与优化，不能默认把它当成基础 RoPE 已经包含的步骤。

\subsection{量化与旋转的顺序通常不可交换}
设 $Q$ 是将实向量映射到有限离散值集合的量化算子，一般有 $Q(\mathcal R(m)\mathbf x)\ne\mathcal R(m)Q(\mathbf x)$。因为量化包含舍入或截断，不是任意旋转下保持不变的线性变换。因而量化缓存中保存旋转前还是旋转后的键，会影响误差结构以及频率切换时能否可靠重建。验证量化系统时应比较注意力分数或最终 logits 的误差，并覆盖不同位置和内容幅度；仅比较单个旋转矩阵是否正交不足以说明量化后精度。

\section{为什么训练长度之外可能失效}
\label{sec:rope-extrapolation}

\subsection{可计算范围、训练范围与有效范围}
给定任意整数位置，式\eqref{eq:rope-block}都可以写出来，这叫公式具有可延伸的定义域。训练时模型只接触某个长度分布，因而只在相应的内容与位置组合上获得优化信号，这叫训练覆盖范围。推理时是否能够从长文本中准确取证、组合信息并生成答案，是有效使用范围。三者不同：运行时接受一个长张量只能证明接口和资源允许，不能证明模型已经学会正确处理其中所有相对距离。

例如训练序列最长为 $L_0$ 时，因果打分中的位移通常落在 $[-L_0+1,0]$。测试长度为 $L_1>L_0$ 后，某些位移超出该区间；即使高频平面已经转过多圈，低频平面的相位仍可能进入训练期间没有覆盖的区域。多个频段的联合相位与内容组合也可能发生变化。将问题概括为“正弦函数见过所有数值所以不需要外推”忽略了联合表示与任务数据分布。

\subsection{外推问题不要求分数无界}
式\eqref{eq:rope-partial}给出了固定内容下的范数界，说明纯旋转不会令单个向量无限膨胀。外推失效却仍然可能发生：一个原本应该较小的错误候选分数在新位移上升高，就可能改变 softmax 排序；新的层间相互作用还会改变后续隐藏状态。失效只需要相对排序与训练目标不相符，并不需要数学意义上的无穷大。讨论异常注意力时，应区分相对于训练分布的过大分数与没有上界的理论发散。

训练时“经常见过某个位置”也不等于“见过该位置上足够多的检索任务”。例如大量样本虽然填满窗口，但答案主要依赖最近几十个词元，模型可能很少获得利用远处证据的有效监督。反过来，针对远距离检索训练，也不能自动保证跨多段证据推理。位置机制决定如何表达结构，数据和损失决定哪些结构会被学会使用。

\subsection{长度增加还改变 softmax 的竞争环境}
设一个正确证据的分数为 $a$，其余 $N-1$ 个干扰候选的分数均为 $b_0$，则正确证据的权重为
\begin{equation}
p_{\mathrm{hit}}=\frac{1}{1+(N-1)\exp(b_0-a)}.
\label{eq:rope-distractor}
\end{equation}
即使相关证据与查询的相对位置完全不变，单纯增加干扰数也会降低它的归一化权重。若希望这个权重至少为 $p\in(0,1)$，需要 $a-b_0\geq\log(N-1)+\log(p/(1-p))$。这是一个简化模型，真实分数并不相同，但它说明长度问题还包括候选竞争，而不只是频率表覆盖。

因此，扩展位置编码后只观察训练损失下降，不能断言远处信息已经被可靠利用。应该固定证据内容，分别改变距离、干扰数量、证据所在位置和问题难度，才能分辨性能变化来自位置几何、注意力竞争还是任务本身。有关注意力温度的修改正是在这类背景下具有意义，但其有效性仍需经验校准。

\section{长上下文扩展：位置、频率和分数的不同修改}
\label{sec:rope-scaling}

\subsection{统一写法与静态方案的平移条件}
为了比较扩展方法，可把第 $j$ 个平面的相位一般地写为 $\phi_j(m)$。相同构造给出相对角度 $\phi_j(n)-\phi_j(m)$。如果 $\phi_j(m)=a_jm+c_j$，共同平移 $m,n$ 后差仍只依赖 $n-m$，常数 $c_j$ 在查询、键采用相同约定时抵消；如果 $\phi_j$ 是非线性位置映射，则这个差一般还依赖两个位置所在的区域。因而“仍然使用旋转”与“仍然只显式依赖相对位移”并非同义。

扩展方法还可能改变分数尺度或分段选择参数。比较它们时，至少应明确三件事：位置坐标是否压缩，哪些频率改变，注意力内积是否附加了倍率。模型配置中的一个字符串只标识某个实现家族，具体公式仍由版本、检查点和配置参数决定。Hugging Face 的 RoPE 工具文档列出了多种类型及各自字段\cite{hf-rope}；迁移配置时应核对实际代码，不应把不同类型中的同名 \texttt{factor} 无条件解释为完全相同的数学量。

\subsection{位置插值及其局部代价}
\keyterm{位置插值（position interpolation, PI）}的一种基本形式是取扩展因子 $s>1$，把位置 $m$ 映射为 $m/s$，等价于将全部频率变为 $\omega_j/s$。它使较长序列占据与原训练区间相近的相位范围，原始工作将这一变换与继续训练结合\cite{chen2023}。下面只分析这个线性变换本身，不将论文中的模型效果推广为任意检查点都具有的保证。

若希望将新序列端点 $L_1-1$ 精确映射到旧端点 $L_0-1$，则可取 $s=(L_1-1)/(L_0-1)$，其中 $L_0>1$；工程上也常按窗口长度比 $L_1/L_0$ 定义缩放因子。两种约定近似相同却不完全一致，应与配置和训练过程配套。实数位置并无数学障碍，因为旋转矩阵本来就对实相位有定义，不需要先把插值后的索引取整；取整反而会把多个不同位置强制合并。

线性插值的代价可以由式\eqref{eq:rope-chord}直接看出。原来相邻位置相差 $\omega_j$ 弧度，插值后只差 $\omega_j/s$；小角度区间内的表示差也约缩小到原来的 $1/s$。例如把窗口扩展四倍时，所有频段的局部位置差都同时被压缩，查询、键投影原先学习的局部匹配尺度可能不再合适。是否能通过继续训练重新适应，是模型与数据的问题，不能由“相位不超界”单独证明。

均匀插值仍保持严格的相对位置形式，因为 $n/s-m/s=(n-m)/s$。这有利于复用原有旋转实现，但不表示扩展后的模型与原模型在短输入上输出相同：同一个非零距离的相位已被改变。只要 $s\ne1$，除去特殊向量或特殊角度，分数通常就会变化。因此短上下文回归测试应与长上下文测试一起进行。

\subsection{改变基数：推导一种非均匀缩放}
另一种思路是保留最高频率，同时更多地压缩低频。给定 $d_r>2$，若希望最慢频率正好变为原来的 $1/s$，而仍采用几何频率，就应要求
\[
(b')^{-(d_r-2)/d_r}=b^{-(d_r-2)/d_r}/s.
\]
解得
\begin{equation}
b'=b\,s^{d_r/(d_r-2)},\qquad
\omega'_j=(b')^{-2j/d_r}
=\omega_j s^{-2j/(d_r-2)}.
\label{eq:rope-base-change}
\end{equation}
这条代数推导清楚地显示：$j=0$ 的频率不变，最后一个频率缩小 $s$ 倍，中间频率按指数渐变。它解释了一类常被称为 NTK-aware 缩放的公式结构；相关命名和方法讨论可参见文献\cite{peng2023}。本文这里只证明频率端点关系，并不将名称当作已证明完整 Transformer 神经正切核保持不变的依据。

当 $d_r=2$ 时，频率表只有 $\omega_0=1$，改变 $b$ 无法改变该频率，式\eqref{eq:rope-base-change}的分母为零正好暴露了构造不适用的情形。即使 $d_r>2$，这个 $s$ 也只是频率设计参数，不能仅凭公式保证任务有效长度增长 $s$ 倍。高频没有压缩，就仍在更长位移上继续演化；中间频段也没有全部被严格映射到原训练相位范围。因此改变基数与均匀插值承担了不同的局部结构和分布变化权衡。

\subsection{分频插值与注意力温度}
可以为每个频段指定 $\gamma_j\in[0,1]$，构造
\begin{equation}
\omega'_j=(1-\gamma_j)\frac{\omega_j}{s}+\gamma_j\omega_j.
\label{eq:rope-frequency-blend}
\end{equation}
当 $\gamma_j=0$ 时完全插值，当 $\gamma_j=1$ 时保留原频率。通过按训练窗口内的旋转圈数 $L_0\omega_j/(2\pi)$ 选择过渡区，可以在保留局部分辨率与限制低频相位变化之间折中。YaRN 将分频处理与注意力尺度调整结合\cite{peng2023}；这里的式子用于解释分频混合的机制，不能代替特定库中关于过渡索引、取整和倍率的完整实现。

若注意力分数改为 $s_{mn}/T$，其中 $T>0$ 称为\keyterm{温度（temperature）}，小温度使已有分数差被放大，大温度使分布更平缓。若通过把旋转后查询、键都乘以同一个正数 $\alpha$ 实现，则点积会乘以 $\alpha^2$，等价于 $T=1/\alpha^2$。把“向量放大两倍”误解为“分数放大两倍”会得到错误配置；只放大一侧才使点积乘以 $\alpha$。

温度不会改变分数的大小排序，却会改变权重集中程度。对于式\eqref{eq:rope-distractor}中的正确排序，适当放大分数差可能减轻大量干扰的归一化稀释；若错误候选原本排在第一位，放大差距也可能强化错误。因而温度调整不是检索正确性的证明，也无法替代对频率和训练分布的适配。某个经验倍率只能作为指定实验条件下的选择，不能当成普适常数。

\subsection{非均匀搜索、动态方案与训练适配}
更一般的设计可为每个频率配置缩放因子 $s_j$，甚至按位置区间采用不同映射。LongRoPE 原始工作利用非均匀插值搜索、逐步扩展和短上下文适配来构建更长窗口\cite{ding2024}。这些方法说明扩展问题可以被视为对位置变换族的选择与训练适配问题，但论文给出的窗口长度仍属于其模型、数据和评估条件，不能据此推断任意使用相同方法名的模型都具有相同能力。

从数学上看，固定的逐频率线性缩放仍满足相对位置恒等式；按位置分段的非线性映射可能只在各区间内部具有简化形式，跨边界时需要计算真实相位差。按当前序列长度动态选择频率，则涉及第\ref{sec:rope-cache}节的缓存一致性。设计选择越灵活，越需要把参数选择规则、切换时点和缓存语义写清楚，否则相同输入可能因分块方式不同而得到不同结果。

继续训练时，应同时提供足够的长距离监督与原有短任务覆盖。单纯把短样本拼接得更长，可能增加词元数量却没有增加需要长距离依赖的训练信号；只训练长检索任务，也可能损害原有局部能力。评价应覆盖短输入回归、长证据检索、跨段组合和自然文本建模，并记录训练预算。RoPE 扩展是在位置几何、模型适配与系统资源之间协调，不是修改配置后便自动完成的一步。

\section{与其他位置机制的比较}
\label{sec:rope-comparison}

\subsection{可学习绝对位置与正弦位置向量}
可学习的绝对位置嵌入为每个位置配置一组参数，优点是结构直接，模型可以学习位置特有的模式；超过嵌入表覆盖范围时，则需要额外定义新位置如何获得参数。固定正弦位置向量把位置映射为不同频率的正弦、余弦，能够在公式上生成更长位置的向量。对于未经任意投影的纯位置向量，适当配对的正弦、余弦内积可以化为余弦位移项；但加入内容、线性投影与层间非线性以后，不能自动保留这一纯位置内积性质。

RoPE 与加法正弦编码都使用谐波函数，却有不同的接口。加法方式为输入表示增加位置分量，RoPE 则在注意力头内按位置改变内容向量的方向。前者可以让后续网络学习如何组合内容与位置，后者通过查询、键的配套变换预先确保局部打分具有相对位移形式。实际任务中哪一种效果更好，还取决于架构、训练与预算；仅凭一个形式更紧凑或一个恒等式更漂亮，无法完成经验比较。

\subsection{相对位置偏置与内容调制的差别}
一种相对位置机制直接向分数加入标量偏置 $\beta_{n-m}$，得到 $s_{mn}=\mathbf q_m^{\mathsf T}\mathbf k_n/\sqrt d+\beta_{n-m}$。偏置可以是精确距离表、分桶参数或某个函数，它独立于当前查询、键内容。RoPE 则得到 $\mathbf q_m^{\mathsf T}\mathcal R(n-m)\mathbf k_n/\sqrt d$，距离通过矩阵作用与内容相乘。因此两个同距离词元对在偏置项上接受相同加法调整，在 RoPE 中却可能产生不同符号、不同幅度的位置效应。

以因果注意力为例，ALiBi 使用按头设置的线性距离惩罚，其基本形式可写为 $s_{mn}=\mathbf q_m^{\mathsf T}\mathbf k_n/\sqrt d-\eta_h(m-n)$，其中 $\eta_h>0$ 是第 $h$ 个头的斜率，$n\leq m$\cite{press2021}。这一偏置项本身随历史距离增大而降低，但加上内容分数后的总分数仍不必单调。它与 RoPE 的振荡式内容调制有不同的归纳偏置，也再次说明必须区分某个位置项的性质与整个注意力分数的性质。

\subsection{位置机制与高效注意力解决不同问题}
更快的注意力内核可以通过分块与重计算减少显存访问或中间矩阵存储，而不必改变注意力的数学结果。RoPE 规定查询、键在打分前怎样变换，内核规定这些变换与矩阵运算怎样执行；只要接口和数值约定相容，二者可以组合。由此不能推出 RoPE 会降低稠密注意力的成对计算数量，也不能把某个融合内核的吞吐量提升归因于位置编码本身。

对于使用特征映射的线性注意力，情况需要重新分析。设特征向量为 $\varphi(\mathbf q),\varphi(\mathbf k)$，则“先旋转原查询再做非线性映射”与“先映射再旋转特征”一般不等价。若依赖非负特征保证归一化分母为正，旋转还可能引入负坐标，破坏原来的论证。RoPE 的双线性相对位置恒等式仍可作为构造工具，但标准 softmax 注意力的推导不能原封不动转移到任意核化或线性注意力架构。

\section{进阶视角：群结构、多轴位置与信息解释}
\label{sec:rope-advanced}

\subsection{从位置加法到旋转复合}
整数位置在加法下构成群，而 $m\mapsto\mathcal R(m)$ 把位置相加映射为正交矩阵相乘。这可以称为一个\keyterm{正交表示（orthogonal representation）}。不需要群论背景也能理解它的用途：连续移动两次与一次移动总位移产生相同变换，逆向移动对应逆矩阵。式\eqref{eq:rope-main}之所以简洁，正是因为查询上的逆向变换与键上的正向变换合并为相对移动。

这个视角还帮助区分“可学习”与“结构保持”。可以学习频率，只要对查询与键使用同一组固定后的频率，位置加法规律仍成立；也可以用固定正交矩阵 $\mathbf U$ 变换所有坐标平面，构造 $\mathbf U\mathcal R(m)\mathbf U^{\mathsf T}$，相对恒等式仍成立。若独立学习每个位置的任意正交矩阵，虽然每个矩阵仍保持范数，却通常不再满足组合规律。正交性保证局部长度，组合规律保证相对位移，两种条件不可互相替代。

\subsection{推广到二维图像与多轴坐标}
图像、视频和多模态输入可能具有多个位置轴。设位置为 $\mathbf p\in\mathbb R^a$，例如 $a=2$ 表示图像的行列坐标，给第 $j$ 个旋转平面指定频率向量 $\boldsymbol\omega_j\in\mathbb R^a$，使用相位 $\phi_j(\mathbf p)=\boldsymbol\omega_j^{\mathsf T}\mathbf p$。则查询位置 $\mathbf p$ 与键位置 $\mathbf u$ 的相对角度为 $\boldsymbol\omega_j^{\mathsf T}(\mathbf u-\mathbf p)$，相同代数推导仍然成立。

一种简单设计把部分平面分给行轴、其他平面分给列轴，即频率向量只在一个轴上非零；另一种设计允许每个频率同时混合多个轴。两者在表达假设、参数组织和训练适配上不同，本文给出的线性相位构造只是数学例子，不意味着所有视觉模型都按同一规则实现。文本顺序、图像网格和视频时间也具有不同单位，合并前必须明确坐标尺度与编号规则。

尤其要避免“二维 RoPE”的歧义：基础文本 RoPE 已经把特征拆成二维实平面，而视觉中的二维位置通常指行、列两个位置轴。这两个“二维”位于不同对象上，一个是特征空间中的旋转平面，另一个是输入的空间坐标。即使位置有三个轴，每个基本旋转块仍可只有两个实坐标，通过一个由三个位置分量共同决定的相位连接起来。

\subsection{为什么不能从一个旋转向量直接读出位置}
若只观察 $\widetilde{\mathbf q}_m=\mathcal R(m)\mathbf q_m$，而不知道未旋转的 $\mathbf q_m$，则一般无法唯一恢复 $m$。对任意候选位置 $t$，都可以定义 $\mathbf q_t'=\mathcal R(t)^{\mathsf T}\widetilde{\mathbf q}_m$，使 $\mathcal R(t)\mathbf q_t'=\widetilde{\mathbf q}_m$。因此位置不是作为一个可直接拆下的标签附在向量末尾，而是通过共同约定影响内容之间的匹配关系。学习到的内容分布可能使位置在统计上可预测，但那是额外的分布信息。

类似地，如果通过线性探针从隐藏状态中预测出词元位置，也不能仅据此断言某一频段“专门存储绝对位置”。起始符、上下文长度、句子结构和训练数据都可能提供相关线索。要研究某个频段的因果作用，应在控制内容和其他通道的条件下进行干预，例如改变该频段相位、置零相应查询键坐标，再比较特定任务的变化。相关性分析、代数可识别性和因果贡献属于不同证据层次。

\section{从公式正确到系统可靠的验证方案}
\label{sec:rope-validation}

\subsection{用结构性质检查局部实现}
局部检查首先固定双精度随机向量、频率、位置及坐标布局。位置为零时应返回原旋转部分；旋转后范数应在容差内保持；分别旋转查询和键得到的内积应等于相对旋转形式；同时给两者位置增加常数，应保持点积。这几项覆盖了角度、转置和位置符号的主要错误。只检查范数远远不够，因为错误的频率、错误的配对乃至反向旋转也可能保持范数。

测试不宜只用全一向量或相同查询键，因为对称输入可能掩盖正弦项错误。应包含正负混合、不同幅度、仅一个平面非零以及未旋转尾部非零的向量，并用第\ref{sec:rope-rotation}节的正交基算例直接检查方向。形状应覆盖不同批量、不同头数、不同查询键长度和部分旋转。相邻、半段与复数实现的比较必须先明确坐标置换，才能把不一致解释为真正的实现错误。

梯度检查可在小维数双精度输入上，用中心有限差分比较标量损失对查询坐标的导数与式\eqref{eq:rope-backward}。扰动步长不能过小以致舍入误差主导，也不能过大以致截断误差明显。对频率、位置或缩放因子求导时，要确认它们确实被设计为可微输入；离散词元索引通常只是外部数据，不需要在训练中对其求导。

\subsection{用完整路径检查缓存语义}
固定一小段输入，关闭 dropout，比较整段因果前向与逐词缓存前向的每个位置输出；再改变预填充分块大小，确认同一逻辑序列的结果在容差内一致。这比只验证最后一个词元更容易定位错误。如果第一块一致、第二块开始偏离，应检查位置偏移、缓存写入索引和因果掩码；如果只在不同填充方式下偏离，应检查批次内每个样本的位置与有效长度。

缓存测试还应包含非零起点、窗口裁剪、重复使用同一前缀以及查询头数不同于键值头数的情况。若实现动态频率切换，应分别定义两种参照：一是固定隐藏向量下的缓存旋转转换，二是整个前缀在新配置下重新计算。两者预期不同，不可用前者通过来宣称后者也已等价。量化缓存则应额外测量误差随重复转换次数和序列长度的变化。

\subsection{长度泛化实验的受控变量}
评估长上下文时，可以构造具有确定答案的键值检索任务：在文本若干位置插入随机生成的唯一键及其值，在末尾询问其中一个键。随机键值减少模型依靠语义常识猜测答案的机会，程序能够直接生成真值。随后分别改变总长度、证据到问题的距离、相似干扰项数量以及证据所在的相对位置。每次只改变一个主要因素，更有助于把位置失效与内容混淆区分开。

单条“针在草堆”任务通过，只能证明在该协议下能够找回一条证据。还应增加需要连接两条远隔记录的组合任务、需要比较顺序的任务以及带冲突干扰的任务。自然文本评价可补充语言建模损失与摘要或问答质量，但要控制截断方式、问题长度和答案评价规则。测试答案若在提示中重复出现，或基线可通过最近位置的泄漏线索回答，就无法证明模型真正读取了远处目标。

公平比较原始 RoPE 与扩展方案时，需要记录原检查点、训练长度分布、额外训练词元数、优化配置、精度和推理窗口。若一个方法经过额外继续训练，另一个只修改推理配置，其差异同时包含训练预算，不能全部归因于位置公式。超参数只能在验证集上选择，最终长度与任务测试应留出独立样本。对于随机数据生成，保存随机种子与生成协议，使失败样例可以重新定位。

\subsection{评价成本与失败，而不只报告最大窗口}
有效窗口不是一个脱离任务的单一数字。可以为每个长度报告准确率、证据位置分布、每类错误比例与置信区间，同时报告预填充时间、逐词解码时间和缓存占用。即使任务质量不变，窗口增大也可能显著改变吞吐量；即使系统能装入更长上下文，检索质量也可能先下降。将任务性能和资源消耗放在同一个比较协议中，才能判断扩展是否具有实际价值。

出现长输入失败时，优先检查可确定的实现问题，再分析模型能力。位置索引、旋转次数、坐标排列、掩码和缓存版本都可通过小样例精确验证；这些通过以后，才更有依据研究训练覆盖、频率选择和远距离监督是否不足。本文的验证方案是可执行的实验设计，不预设某种缩放方法必然胜出，也不把尚未运行的实验写成已有结果。

\section{习题与详细解答}
\label{sec:rope-exercises}

\subsection{相同旋转为什么不能提供相对距离}
\begin{example}[固定角度与位置角度]
有人提出在所有位置给查询和键统一旋转 $\pi/3$，认为这样也使用了旋转位置编码。证明这一变换是否改变注意力分数，并说明与 RoPE 的区别。
\end{example}
\begin{solve}
设共同正交矩阵为 $\mathbf U$，则 $(\mathbf U\mathbf q_m)^{\mathsf T}(\mathbf U\mathbf k_n)=\mathbf q_m^{\mathsf T}\mathbf U^{\mathsf T}\mathbf U\mathbf k_n=\mathbf q_m^{\mathsf T}\mathbf k_n$。所以在只通过内积使用查询、键的条件下，所有分数不变，也就没有新增距离信息。RoPE 给不同位置分配一般不同的角度，使两个旋转合成后保留 $n-m$。关键在于变换与位置的关系，而不是是否出现三角函数。
\end{solve}

\subsection{方向信息与正弦项}
\begin{example}[前后位置可否区分]
在一个频率为 $\omega=\pi/4$ 的平面中，取 $\mathbf q=(1,0)^{\mathsf T}$、$\mathbf k=(0,1)^{\mathsf T}$。分别计算 $\delta=1$ 和 $\delta=-1$ 的未缩放分数；再把键改为 $(1,0)^{\mathsf T}$，比较结果。
\end{example}
\begin{solve}
第一种内容对满足 $A=0,B=-1$，因此分数为 $-\sin(\delta\pi/4)$。两个位移分别得到 $-\sqrt2/2$ 和 $\sqrt2/2$，可以区分方向。第二种内容对满足 $A=1,B=0$，分数为 $\cos(\delta\pi/4)$，两个方向都得到 $\sqrt2/2$。这说明方向敏感性取决于内容系数与频段组合；RoPE 允许模型表达方向，不强迫每个查询键对都具有非对称分数。
\end{solve}

\subsection{错误坐标配对为什么不会被范数测试发现}
\begin{example}[四维置换]
取向量 $\mathbf x=(1,2,3,4)^{\mathsf T}$。写出与相邻配对对应的半段排列，并说明直接将原数组送进半段旋转函数为何可能通过范数检查却改变注意力。
\end{example}
\begin{solve}
相邻配对为 $(1,2)$、$(3,4)$，所以半段排列应为 $(1,3,2,4)$，其前后半段重新配对后仍得到这两个平面。若直接输入 $(1,2,3,4)$，半段函数会旋转 $(1,3)$、$(2,4)$，混合的是不同特征。两种旋转矩阵都正交，因此都保持 $1^2+2^2+3^2+4^2=30$，范数测试无法区分。应在正确置换下比较输出，或用明确的基向量检查每个频率作用于哪些坐标。
\end{solve}

\subsection{缓存键究竟应该旋转几次}
\begin{example}[重复旋转的相位]
位置 $n$ 的键已缓存为 $\widetilde{\mathbf k}_n=\mathcal R(n)\mathbf k_n$，下一步程序又对整个缓存按原位置旋转一次。写出查询位置 $m$ 得到的内积，并说明它为什么错误。
\end{example}
\begin{solve}
缓存键会变成 $\mathcal R(n)^2\mathbf k_n=\mathcal R(2n)\mathbf k_n$。查询仍为 $\mathcal R(m)\mathbf q_m$，所以内积为 $\mathbf q_m^{\mathsf T}\mathcal R(2n-m)\mathbf k_n$，目标却是 $\mathbf q_m^{\mathsf T}\mathcal R(n-m)\mathbf k_n$。除去特殊向量或相位重合，二者不同。正确实现应明确缓存保存旋转前还是旋转后键；保存后者时，固定频率下旧键无需每步重新旋转。
\end{solve}

\subsection{只缩放一侧为何破坏共同平移性质}
\begin{example}[查询与键使用不同位置尺度]
设查询采用 $\mathcal R(m/s)$，键仍采用 $\mathcal R(n)$，其中 $s>1$。推导内积中的位置变量，并计算共同平移 $c$ 后多出的项。
\end{example}
\begin{solve}
内积中出现 $\mathcal R(n-m/s)$，这不是仅由 $n-m$ 决定的函数。平移后位置变量变为 $(n+c)-(m+c)/s=n-m/s+c(1-1/s)$，多出与原点有关的项。因此位置插值必须对参与点积的两侧采用相容规则。缓存中若仍存旧尺度的键，而查询已经使用新尺度，就会出现同类不一致；这不应被解释为正常的长度泛化误差。
\end{solve}

\subsection{由相对位置推导缓存重编号}
\begin{example}[窗口起点移动]
一个窗口当前保存原位置 $100$ 至 $199$ 的旋转后键，希望把它们重新编号为 $0$ 至 $99$。给出转换公式，并说明为什么不能把这一步等同于截断输入后重新运行模型。
\end{example}
\begin{solve}
对所有保存键施加共同旋转 $\mathcal R(-100)$，则 $\mathcal R(-100)\mathcal R(n)\mathbf k_n=\mathcal R(n-100)\mathbf k_n$。查询也使用减去 $100$ 的新位置，保留内容间的相对内积不变。然而原隐藏状态可能曾经读取位置 $0$ 至 $99$ 的内容；截断后重新运行会改变这些隐藏状态，进而改变查询、键和值。重编号只改变坐标表示，并没有抹去历史计算已经汇入的内容。
\end{solve}

\subsection{相位范围缩小不等于注意力不变}
\begin{example}[插值对短距离的影响]
在单个平面中取 $\mathbf q=\mathbf k=(1,0)^{\mathsf T}$、$\omega=1$，比较距离 $\delta=1$ 在原频率和四倍线性插值后的分数。它是否证明插值总能提高局部注意力？
\end{example}
\begin{solve}
原分数为 $\cos1\approx0.540302$，插值后为 $\cos(1/4)\approx0.968912$，在这个特定内容方向上确实升高。但若取 $\mathbf k=(-1,0)^{\mathsf T}$，两者分别为 $-\cos1$ 与 $-\cos(1/4)$，分数反而降低；若取其他方向，还会出现正弦项。并且注意力权重取决于所有候选分数。因此一个同向向量例子不能推出对任意内容的单调改进，更不能推出任务准确率必然提高。
\end{solve}

\subsection{干扰项数量与分数差}
\begin{example}[扩窗后的权重稀释]
在式\eqref{eq:rope-distractor}的简化模型中，若希望正确位置权重至少为 $1/2$，从 $N=100$ 增加到 $N=10000$ 后，所需最小分数差增加多少？这一现象是否完全由 RoPE 造成？
\end{example}
\begin{solve}
目标 $p=1/2$ 使 $\log(p/(1-p))=0$，所需分数差分别为 $\log99$ 与 $\log9999$，增量为 $\log(9999/99)=\log101\approx4.61512$。这里把全部干扰分数设为相同只是分析用假设，但结论清楚显示归一化分母会随候选数改变。推导没有使用旋转公式，所以它不是 RoPE 独有的问题；位置扩展和注意力竞争应分别分析。
\end{solve}

\section{综合实践与进一步阅读}
\label{sec:rope-practice}

\subsection{一个可以分阶段完成的小型实践}
第一阶段不训练网络，只实现二维旋转与高维配对，重现第\ref{sec:rope-example}节的分数，验证相对恒等式和坐标置换。第二阶段构造一个小型因果注意力模块，显式传入查询、键位置，实现旋转后键缓存，比较完整前向、分块前向和逐步前向。第三阶段再接入词嵌入、输出头和训练循环，用程序生成的复制、检索或顺序任务检查网络是否利用了位置结构。每一步都有独立目标，避免在尚未确认底层约定时用大规模训练掩盖错误。

设计对照时，可以比较无显式位置编码、加法正弦编码、基础 RoPE 与一种明确写出公式的缩放方案。保持模型参数预算、训练样本和优化过程尽可能一致，同时分别报告训练范围内与范围外的表现。由于无位置模型仍可能从因果掩码获得部分顺序线索，不能预设其准确率一定等于随机猜测；应实际测量，并用失败样例分析任务是否真的要求精细距离。对照实验的价值在于解释机制，而不是确保某种方法获胜。

所有结果应注明实现检验与模型训练的边界。数值恒等式通过说明旋转公式实现正确，不能说明一个语言模型已经学会长程推理；合成任务通过说明在特定生成分布下有相应能力，也不能直接替代自然文本评价。最有用的实践记录是同时保留配置、输入位置、关键中间张量和错误样例，使一个分数异常能够回溯到具体内容、距离和频段，而不是只留下一个总体准确率。

\subsection{阅读原始文献时应核对什么}
阅读 RoPE 相关材料，首先统一其位移方向、向量行列约定、复共轭位置和坐标配对，再比较公式。其次区分作者证明的代数性质、用于解释的统计假设以及实验支持的结论范围。最后检查实现中是否包含部分旋转、分数倍率、动态频率或缓存转换；论文中的简洁公式与库中的完整路径之间，常有必须明确的工程约定。本文基础推导覆盖的是这些约定共同依赖的核心数学结构。

参考文献中，文献\cite{su2021}用于定位 RoPE 的原始方法，文献\cite{vaswani2017,press2021}用于理解不同位置机制的背景，文献\cite{chen2023,peng2023,ding2024}用于了解几类长度扩展方案。代码与配置阅读可参照文献\cite{meta-llama,hf-rope}。这些在线实现可能随版本变化；本文给出的自包含定义、反例与算例不依赖某个库字段永久不变。复现实验时应保存具体提交或版本及完整配置，再对照本文的数学契约逐项核验。

\begin{thebibliography}{99}
\bibitem{su2021}
Su, J., et al.
\emph{RoFormer: Enhanced Transformer with Rotary Position Embedding}.
arXiv:2104.09864, 2021; revised version, 2023.
\url{https://arxiv.org/abs/2104.09864}.

\bibitem{vaswani2017}
Vaswani, A., et al.
\emph{Attention Is All You Need}.
Advances in Neural Information Processing Systems, 2017.
\url{https://arxiv.org/abs/1706.03762}.

\bibitem{press2021}
Press, O., Smith, N. A., and Lewis, M.
\emph{Train Short, Test Long: Attention with Linear Biases Enables Input Length Extrapolation}.
arXiv:2108.12409, 2021; ICLR, 2022.
\url{https://arxiv.org/abs/2108.12409}.

\bibitem{chen2023}
Chen, S., Wong, S., Chen, L., and Tian, Y.
\emph{Extending Context Window of Large Language Models via Positional Interpolation}.
arXiv:2306.15595, 2023.
\url{https://arxiv.org/abs/2306.15595}.

\bibitem{peng2023}
Peng, B., et al.
\emph{YaRN: Efficient Context Window Extension of Large Language Models}.
arXiv:2309.00071, 2023; ICLR, 2024.
\url{https://arxiv.org/abs/2309.00071}.

\bibitem{ding2024}
Ding, Y., et al.
\emph{LongRoPE: Extending LLM Context Window Beyond 2 Million Tokens}.
arXiv:2402.13753, 2024.
\url{https://arxiv.org/abs/2402.13753}.

\bibitem{meta-llama}
Meta.
\emph{Llama reference implementation: llama/model.py}.
\url{https://github.com/meta-llama/llama/blob/main/llama/model.py}.
访问日期：2026 年 10 月 3 日。

\bibitem{hf-rope}
Hugging Face.
\emph{Transformers documentation: Rotary embeddings utilities}.
\url{https://huggingface.co/docs/transformers/main/en/internal/rope_utils}.
访问日期：2026 年 10 月 3 日。
\end{thebibliography}
\end{document}
